% sections/6_conclusions.tex
\chapter{Conclusioni}
\label{chap:conclusions}

Questo capitolo riassume il lavoro svolto e risponde ai quesiti di ricerca posti nella
Sezione~\ref{sec:intro_questions}. Discute poi i limiti emersi e i possibili sviluppi futuri.

% ---------------------------------------------------------------------------------------------
\section[Sintesi del lavoro]{Sintesi del lavoro}\label{sec:concl_summary}

La tesi ha re-implementato in Python il framework TREAD di Pallotta et al.~\cite{pallotta} e lo ha
applicato ai messaggi AIS pubblicati dal NOAA per l'arcipelago delle Hawaii nel primo trimestre del
2019. La pipeline si articola in tre fasi:
\begin{itemize}
    \item il pre-processing, il cui compito è quello di pulire i dati e classificare gli stati delle navi;
    \item la fase di Knowledge Discovery, che estrae senza supervisione i waypoint, tramite H-DBSCAN, e le rotte che li collegano;
    \item la fase di anomaly detection, che assegna a ogni trip di una rotta un punteggio di normalità, costruito come log-probabilità a partire da due KDE e da una distribuzione di Weibull.
\end{itemize}

Il metodo è stato valutato sulla rotta dal porto di Honolulu al porto di Kalaeloa, sull'isola di
Oahu, confrontando il ranking dei 184 trip con un'etichettatura manuale usata esclusivamente come
riferimento per il calcolo delle metriche.

% ---------------------------------------------------------------------------------------------
\section[Risposte ai quesiti di ricerca]{Risposte ai quesiti di ricerca}\label{sec:concl_questions}

\subsection[Applicabilità del framework a un dataset reale]{Applicabilità del framework a un dataset reale}

Il primo quesito chiedeva se TREAD, sviluppato agli inizi degli anni 2010, sia ancora applicabile
con gli strumenti attuali a un dataset AIS reale di milioni di record. La risposta è positiva,
a condizione di alcuni adattamenti. La pipeline ha elaborato 5.462.415 record, ridotti a 5.415.129
dopo la pulizia, su un computer personale e senza infrastrutture dedicate.

Gli adattamenti hanno riguardato soprattutto la gestione della memoria. L'elaborazione avviene in
modalità batch, a posteriori sull'intero dataset, invece che messaggio per messaggio come
nell'algoritmo originale. Prima del clustering le coordinate dei punti stazionari vengono
arrotondate e aggregate, con un peso logaritmico che tiene conto di quante volte ogni posizione
compare nel dataset. Senza questa aggregazione i circa 3.770.000 punti stazionari non sarebbero
stati gestibili con la memoria disponibile. L'importazione dei file giornalieri in parallelo con
Polars e l'uso in Pandas di tipi di dato compatti, scelti in base alla codifica dei messaggi AIS,
hanno contribuito a contenere l'occupazione di memoria. La sola conversione dei tipi riduce quella
del dataset di circa il 34\%.

Con questi adattamenti il framework produce risultati coerenti con la realtà. I 144 waypoint
estratti comprendono i principali porti dell'arcipelago, le 766 rotte riflettono servizi reali come
il collegamento tra Lahaina e l'isola di Lanai, e sulla rotta selezionata il punteggio di normalità
individua come più anomali trip che deviano in modo evidente dal corridoio.

\subsection[Aspetti riutilizzati e aspetti adattati]{Aspetti riutilizzati e aspetti adattati}

Il secondo quesito chiedeva quali aspetti dell'algoritmo siano riutilizzabili e quali abbiano
richiesto adattamenti o sostituzioni. Dell'impianto originale sono stati mantenuti la suddivisione
in Knowledge Discovery e anomaly detection, la rappresentazione del traffico tramite navi,
waypoint e rotte, la classificazione degli stati delle navi e la struttura del modello
probabilistico. Il punteggio di normalità combina infatti, come in Pallotta et al., un termine di
osservazione, un termine di transizione basato sulla distribuzione di Weibull e la probabilità
a priori della rotta.

Hanno invece richiesto adattamenti o sostituzioni
\begin{itemize}
    \item l'elaborazione, passata da incrementale a batch per i limiti dell'hardware disponibile;
    \item il clustering dei waypoint, eseguito con H-DBSCAN invece che con DBSCAN incrementale;
    \item il termine di osservazione, calcolato su un intorno di 500~m di ogni punto con due KDE stimate una sola volta sull'intera rotta e senza la direzione del moto;
    \item il termine di transizione, con una predizione della posizione semplificata e intervalli temporali scelti in modo adattivo in base alla durata dei trip;
    \item il denominatore della probabilità a priori, riferito alle sole navi della rotta analizzata;
    \item l'unità di valutazione, che è l'intero trip invece di una finestra temporale mobile.
\end{itemize}
Le singole differenze sono descritte e motivate nelle Sezioni~\ref{sec:meth_algorithm}
e~\ref{sec:meth_anomaly}.

\subsection[Passaggi migliorati con gli strumenti attuali]{Passaggi migliorati con gli strumenti attuali}

Il terzo quesito chiedeva quali passaggi del metodo originale possano essere migliorati con gli
strumenti attuali. Il lavoro ne ha individuati due.

Il primo è l'identificazione dei waypoint. H-DBSCAN è stato pubblicato nel 2013~\cite{campello2013}, lo stesso
anno del lavoro di Pallotta et al., che è il risultato di un'attività svolta probabilmente negli anni precedenti.
È quindi difficile che gli autori potessero conoscerlo e adottarlo. A differenza di DBSCAN,
H-DBSCAN non usa un raggio fisso ma si adatta alla densità locale dei punti. Nelle
Hawaii, dove i porti principali come Honolulu hanno una concentrazione di punti molto più alta dei
porti minori e delle zone di passaggio, questa caratteristica permette di separare aree vicine nelle
zone dense senza frammentare quelle sparse, come mostrano le figure della
Sezione~\ref{sec:res_qualitative}. Il miglioramento è osservato sulle mappe e non è stato misurato
con un indice quantitativo.

Il secondo è il calcolo delle componenti della velocità. Pallotta et al.~le definiscono in modo
implicito, tramite relazioni da cui andrebbero ricavate gestendo l'ambiguità di quadrante
dell'arcotangente. In questo lavoro sono calcolate direttamente da \texttt{SOG} e \texttt{COG}
con seno e coseno, in modo più semplice e senza ambiguità.

Un altro passaggio, il calcolo del punteggio nelle zone portuali, si è invece rivelato il
principale limite del metodo. È discusso nella Sezione~\ref{sec:concl_limits}, mentre i
possibili interventi sono proposti nella Sezione~\ref{sec:concl_future}.

\subsection[Metriche per la valutazione delle prestazioni]{Metriche per la valutazione delle prestazioni}

Il quarto quesito chiedeva quali metriche siano adatte a quantificare le prestazioni
dell'algoritmo. Poiché il metodo non è supervisionato e i dati AIS non contengono un'indicazione
delle anomalie, il primo passo è costruire un riferimento. In questo lavoro i 184 trip della rotta
selezionata sono stati etichettati a mano in tre classi, cioè normali, route deviation e altre
anomalie. La valutazione è stata condotta in due casi, considerando tutte le anomalie (caso a) e le
sole route deviation (caso b), per distinguere il tipo di anomalia per cui il metodo è pensato
dalle altre.

Sono stati scelti poi due tipi di metriche:
\begin{itemize}
    \item Le metriche che dipendono da una soglia, cioè la matrice di confusione e i valori 
    di precision (quota dei trip segnalati che sono davvero anomali), di recall (quota delle
    anomalie rilevate) e di F1 che ne derivano, descrivono il comportamento
    del metodo con la soglia effettivamente usata (in questo caso la soglia è stata fissata 
    con la regola di Tukey senza ricorrere alle etichette). 
    Con questa soglia la precision è pari a 1 in entrambi i casi, quindi nessun trip normale viene segnalato come anomalo.
    La recall è pari a 0,220 nel caso (a) e a 0,429 nel caso (b), e la F1 vale rispettivamente 0,360 e 0,600. L'accuracy è pari a 0,826 e 0,927 ma è poco indicativa, perché i trip normali sono la grande maggioranza.
    \item Le metriche indipendenti dalla soglia, cioè la curva ROC e l'area sotto la curva (AUC), valutano invece la capacità del punteggio di ordinare i trip, a prescindere dalla soglia scelta. L'AUC vale 0,751 nel caso (a) e 0,844 nel caso (b).
\end{itemize}
Il confronto tra le due famiglie di metriche è informativo. Una recall bassa insieme a un'AUC alta indica che il punteggio ordina bene i trip ma la soglia di Tukey è troppo severa, mentre la differenza tra i due casi mostra che il metodo rileva le route deviation molto meglio delle altre anomalie.

% ---------------------------------------------------------------------------------------------
\section[Limiti]{Limiti}\label{sec:concl_limits}

Ai limiti della metodologia già discussi nella Sezione~\ref{sec:meth_limits} se ne aggiungono altri
emersi dai risultati.

Il limite principale riguarda le zone portuali. Davanti ai porti le navi rallentano e i punti si
addensano, per cui il termine di osservazione assume i valori più alti della rotta. Le manovre lente
nella zona di avvicinamento al porto e le deviazioni che avvengono vicino al porto ottengono
quindi punteggi alti anziché bassi, e quattro anomalie compaiono tra i dieci trip più regolari del
ranking.

Un secondo limite deriva dal calcolo del punteggio come media sull'intero trip. Un'anomalia limitata
a un breve tratto, come l'arrivo in un punto diverso dal porto abituale, pesa poco sul risultato.
Per questo motivo nessuno dei 15 trip che terminano il viaggio fuori dal porto di arrivo viene
segnalato come anomalo.

La valutazione, infine, è limitata a una sola rotta e a tre mesi di dati, e si basa su
un'etichettatura svolta da un solo annotatore. Un'etichettatura con più annotatori renderebbe la valutazione più solida. I risultati quantitativi vanno quindi letti come
indicativi del comportamento del metodo su una rotta molto regolare, e non come prestazioni
generalizzabili ad altre rotte o ad altre aree.

% ---------------------------------------------------------------------------------------------
\section[Sviluppi futuri]{Sviluppi futuri}\label{sec:concl_future}

I risultati e i limiti emersi suggeriscono alcuni possibili sviluppi.
\begin{itemize}
    \item \textbf{Parametro di scala.} La retta $\alpha = m \cdot \Delta t$ descrive solo in modo approssimativo i parametri di scala della Weibull stimati sui diversi intervalli (Figura~\ref{fig:res_weibull_regression}). Si potrebbero esplorare relazioni non lineari e stimare i parametri su intervalli temporali coerenti con il campionamento dei dati.
    \item \textbf{Zone portuali.} Per ridurre il vantaggio che l'addensamento dei punti davanti ai porti dà al punteggio si potrebbero, per esempio, stimare modelli separati per le zone portuali e per il resto della rotta.
    \item \textbf{Direzione del moto.} Introdurre la direzione nella KDE, come nel modello originale, permetterebbe di riconoscere i trip che attraversano zone frequentate ma con una direzione incompatibile con quella abituale, come nell'esempio della doppia inversione di rotta riportato da Pallotta et al.
    \item \textbf{Generalizzazione.} Applicare il metodo ad altre rotte, come quella tra Lahaina e l'isola di Lanai, ad altre aree e a periodi più lunghi permetterebbe di verificare se i risultati ottenuti valgono anche in altri contesti.
    \item \textbf{Versione online.} Riportare la pipeline al funzionamento incrementale del framework originale, aggiornando il modello e i punteggi all'arrivo di ogni messaggio, renderebbe il metodo utilizzabile per il monitoraggio in tempo reale.
\end{itemize}

% ---------------------------------------------------------------------------------------------
\clearpage
\section[Considerazioni finali]{Considerazioni finali}\label{sec:concl_final}

Il lavoro ha mostrato che il framework TREAD, a più di dieci anni dalla sua pubblicazione, resta una
base valida per l'analisi del traffico marittimo. Con gli strumenti attuali e alcuni adattamenti
estrae in modo automatico una rappresentazione del traffico coerente con la realtà e permette di
individuare le deviazioni dalle rotte abituali senza falsi allarmi sulla rotta analizzata. La
valutazione quantitativa ne ha messo in luce anche i limiti, in particolare nelle zone portuali, e
indica le direzioni in cui il metodo può essere ulteriormente migliorato.
